\documentclass[11pt]{article}
\usepackage[T1]{fontenc}
\usepackage{lmodern}
\usepackage[margin=1.03in]{geometry}
\usepackage{microtype}
\usepackage{amsmath,amssymb,amsthm,mathtools,booktabs}
\usepackage[hidelinks]{hyperref}
\hypersetup{pdftitle={Sharp third and fourth Taylor coefficients in Hardy spaces}}
\newtheorem{theorem}{Theorem}
\newtheorem{proposition}[theorem]{Proposition}
\newtheorem{lemma}[theorem]{Lemma}
\newtheorem{corollary}[theorem]{Corollary}
\newcommand{\D}{\mathbb D}
\newcommand{\T}{\mathbb T}
\newcommand{\C}{\mathbb C}
\newcommand{\R}{\mathbb R}
\newcommand{\Ent}{\operatorname{Ent}}
\newcommand{\Dif}{\mathcal D}
\newcommand{\coef}[1]{[z^{#1}]}
\allowdisplaybreaks[1]
\title{Sharp third and fourth Taylor coefficients\\in Hardy spaces below \texorpdfstring{$H^1$}{H1}}
\author{}
\date{27 November 2026}
\begin{document}
\maketitle

\begin{abstract}
We determine the sharp third and fourth Taylor coefficient norms in
the classical Hardy spaces \(H^p(\mathbb D)\) for every \(0<p<1\) and
classify all extremizers. A reflected stationarity identity keeps the
polynomial factors of extremizers free of zeros on the unit circle,
so unconstrained coefficient variations apply. In degree four, the
zero-free condition excludes nonreal stationary phases and selects a
unique admissible branch of a sextic. For any degree, the same
stationary framework yields a finite algebraic computation at
algebraic exponents. The coefficient norm is algebraic at rational
\(p\) and transcendental at irrational algebraic \(p\). In particular,
\(C(4,1/2)\) has degree five and is not expressible by radicals.
We also relate
the limit \(p\uparrow1\) to a maximum of entropy over unit-root
polynomials.
\end{abstract}

\section{The sharp coefficient problem}

Write \(m\) for normalized arc measure on \(\T\). For \(p>0\)
the classical Hardy norm or quasi-norm of a holomorphic function on \(\D\)
is
\[
 \|f\|_p=\sup_{0<r<1}
       \left(\int_{\T}|f(rz)|^p\,dm(z)\right)^{1/p}.
\]
It equals the boundary quasi-norm when finite \cite{Duren}. For \(n\ge1\), let
\begin{equation}\label{eq:def}
 C(n,p)=\sup_{\|f\|_p\le1}|\coef n f|,\qquad f\in H^p(\D).
\end{equation}
Zeros are allowed in \eqref{eq:def}. Brevig and Saksman
\cite{BrevigSaksman} posed the explicit computation of these constants
in the range \(0<p<1\). Their Theorems 1 and 2 give \(C(2,p)\) and
\(C(3,2/3)\), respectively, and identify the corresponding equality
cases. The 2026 account of Chen, Huang, Jin and Li \cite{ChenEtAl}
records the known analytic cases and the remaining coefficient
problem. Brevig and S{\o}vik \cite{BrevigSovik} then showed that
extremizers are outer and resolved the structural nonvanishing and
strict-monotonicity conjectures of \cite{BrevigSaksman}. The third
and fourth coefficient norms for the full range of \(p\) are the
targets here.

Our method combines polynomial jet comparison and exclusion of
competing stationary branches, as developed for the zero-free
coefficient problems in VMRC \cite{VMRC248,VMRC249}. The auxiliary
problem here is to exclude zeros \emph{on} the boundary of the
disk. Its solution permits variations in all polynomial
coefficients. The degree-four proof then uses a phase calculation
and the strict bound on the sum of the reciprocal roots.
The resulting extremizers are unique up to rotation of the variable
and multiplication by a unimodular constant. Section~\ref{sec:arithmetic}
gives a separate arithmetic consequence in every degree.

\section{Boundary separation and stationary jets}

\begin{proposition}\label{thm:boundary}
Fix \(n\ge1\), \(0<p<1\), and \(s=2/p\). The norm \(C(n,p)\)
is attained. Every extremizer of quasi-norm one has the form
\[
 f=\eta Q^s,\qquad
 \deg Q=n,\quad \|Q\|_2=1,\quad Q(0)>0,\quad
 Q(z)\ne0\ \text{for }|z|\le1,\quad |\eta|=1 .
\]
The power uses the analytic logarithm of \(Q\) which is real at zero.
\end{proposition}

\begin{proof}
The Hardy pointwise estimates give a normal family on every smaller
disk. A locally uniform limit of a maximizing sequence still has
quasi-norm at most one, by integration on each smaller circle.
The coefficient converges and the maximum is positive, since \(z^n\)
is admissible. Scaling shows every extremizer has quasi-norm one.
Brevig--S{\o}vik's outer theorem \cite[Thm.~1]{BrevigSovik} applies
to the bounded linear functional \(\coef n\). For convenience, its
concavity argument is as follows. Write \(f=IO\), where \(I\) is
inner and \(O\) is outer. If \(I\) is nonconstant, set
\(u=2\Re I-c\) on \(\T\), with
\(c=\int_{\T}2\Re I\,|f|^p\,dm\). Then \(u\) is nonconstant and has
mean zero under \(|f|^p\,dm\). The analytic function
\(g=O(1+I^2-cI)\) belongs to \(H^p\), and \(g=uf\) on \(\T\).
For small real \(t>0\), strict concavity of \(x^p\) gives
\(\|f\pm tg\|_p^p<1\), since \(1\pm tu>0\).
Averaging the two coefficient values contradicts extremality.
Thus \(f\) is outer. Here \(\Re I\) cannot be constant on \(\T\)
unless \(I\) is constant, and the outer boundary modulus is positive
almost everywhere.

By rotations assume \(a_0=f(0)>0\) and \(a_n=\coef n f=C(n,p)>0\).
Vary \(f e^{t g}/\|f e^{t g}\|_p\), first with an analytic polynomial
\(g\) and then with \(ig\). Differentiating the norm is justified
by boundedness of \(g\) on \(\T\). Stationarity gives
\[
 \coef n(fg)=a_n\int_{\T}|f|^p g\,dm .
\]
With \(g=z^k\) for every \(k\ge0\), Fourier uniqueness gives
\[
 h(z)=a_n+\sum_{k=1}^n
       \bigl(a_{n-k}z^{-k}+\bar a_{n-k}z^k\bigr)
     =a_n |f(z)|^p \quad\text{on }\T\text{ a.e.}
\]
The Fej\'er--Riesz theorem factors \(h=a_n |Q|^2\) with
\(Q\) free of zeros in \(\D\), \(Q(0)>0\), \(\|Q\|_2=1\), and
degree exactly \(n\) because \(a_0>0\). The normalized outer function
with boundary modulus \(|Q|^{2/p}\) is \(Q^s\), so \(f=Q^s\).

If \(U\) has degree at most \(d\), write \(U^{*d}(z)=
z^d\overline{U(1/\bar z)}\).
Let \(R=z^n h=a_n QQ^{*n}\), let \(\Dif=z\,d/dz\), and set
\(E=R\Dif\log f=s a_n Q^{*n}\Dif Q\).
The polynomial \(R\) is self-reciprocal of degree \(2n\), and
\(R=f\bmod z^{n+1}\), meaning their Taylor coefficients agree through
order \(n\). Consequently the first \(n+1\) coefficients of \(E\)
are \(k a_k\), \(0\le k\le n\). The product rule and
\((\Dif Q)^{*n}=nQ^{*n}-\Dif(Q^{*n})\) give
\[
 E-E^{*2n}=s(\Dif R-nR).
\]
Set \(\beta=1-1/p=1-s/2\). The first \(n+1\) coefficients of
\((E+E^{*2n})/2\) are \((n-\beta(n-k))a_k\).
Reflecting the remaining coefficients shows that, on \(\T\),
\begin{equation}\label{eq:boundary-identity}
 z^{-n}\frac{E+E^{*2n}}{2}=nh-\beta D_h,
\end{equation}
where \(D_h\) is the outward radial derivative of the harmonic
extension of \(h\), the Fourier multiplier \(|k|\).

Suppose \(Q(\zeta)=0\) for some \(\zeta\in\T\). Then \(E(\zeta)=0\)
by its factored form and \(E^{*2n}(\zeta)=0\) by reflection.
After rotation to \(\zeta=1\), write \(Q(z)=(z-1)B(z)\).
Applying the multiplier to
\(h=a_n(2-z-z^{-1})|B|^2\) at 1 gives
\(D_h(1)=-2a_n\|B\|_2^2<0\), since
\(2|k|-|k+1|-|k-1|\) vanishes except at \(k=0\), where it
equals \(-2\). Equation \eqref{eq:boundary-identity} would now give
\(0=-\beta D_h(\zeta)\), which is impossible. This proves strict
boundary separation.
\end{proof}

There are two useful forms of stationarity. Put \(a=s-1>1\),
\(q_0=Q(0)\), and \(P=Q/q_0\). With \(a_n>0\) as in the proof,
the leading coefficient \(d=\coef n P\) is positive because the
constant term of \(R=a_n QQ^{*n}\) is \(q_0^s>0\).
The congruence \(R=f\bmod z^{n+1}\) gives
\begin{equation}\label{eq:jet}
 P^a\equiv P^{*n}/d\pmod{z^{n+1}},\qquad
 a_n=\frac{q_0^{a-1}}{d},\qquad
 q_0=(1+|[z]P|^2+\cdots+|\coef nP|^2)^{-1/2}.
\end{equation}
When rotating the argument without the \(a_n>0\) normalization,
\(P^a\equiv P^{*n}/\bar d\bmod z^{n+1}\) and
\(\coef n Q^s=q_0^{a-1}/\bar d\). These identities also follow
from \eqref{eq:jet} by rotation.
For polynomial coefficient optimization we instead use the complex
derivatives of \(F(q)=\coef n Q^s\) as a local function of
\(q=(q_0,\ldots,q_n)\). Strict boundary separation is an open
condition on the coefficients. All sufficiently small variations
on the complex \(H^2\) unit sphere are therefore admissible.
Allowing the local analytic branch also for complex \(q_0\) near
its positive value, homogeneity gives
\begin{equation}\label{eq:sphere-lagrange}
 \frac{\partial F}{\partial q_j}=s F\,\overline{q_j},
                  \qquad 0\le j\le n
\end{equation}
at a maximum of \(|F|\). Indeed, apply Lagrange multipliers to
\(|F|^2\) and use \(\sum q_j\,\partial F/\partial q_j=sF\).
Under rotation of the argument, \(F\) may be complex.
\section{The third coefficient}

\begin{theorem}\label{thm:third}
For \(0<p<1\), set \(s=2/p\) and \(a=s-1>1\). Let \(u\) be the
smallest positive root of
\begin{equation}\label{eq:third-cubic}
 P_a(u)=\frac{a-2}{6}u^3+\frac{a(5-a)}4u^2
             +\frac{a(1-2a)}2u-a^2 .
\end{equation}
Set
\[
 y=\left(a^2+\frac{a(a-1)}2u\right)^{-1/2},\qquad x=uy,
 \qquad q_0=\left(1+x^2+a^2x^2y^2+y^2\right)^{-1/2}.
\]
Then
\[
 C(3,p)=\frac{q_0^{s-2}}{y}.
\]
The unit extremizers are precisely the unimodular scalar and variable rotations of
\[
 Q(z)^s,\qquad Q(z)=q_0(1+xz+a x y z^2+y z^3).
\]
The power is defined using the analytic logarithm real at zero.
\end{theorem}

\begin{proof}
Proposition~\ref{thm:boundary} represents every extremizer as
\(Q^s\), up to phase, with \(Q\) zero-free on the closed disk.
The open coefficient variations used below are therefore admissible.

Write \(Q(z)=\sum_{j=0}^3 q_j z^j\) and let \(M=C(3,p)>0\). The third
coefficient of \(Q^s\) is
\begin{equation}\label{eq:third-coefficient}
 F=sq_0^{s-1}q_3
     +saq_0^{s-2}q_1q_2
     +\frac{sa(a-1)}6q_0^{s-3}q_1^3 .
\end{equation}
The phases of the three nonzero summands can be varied independently,
since the angular coordinates \(\theta_3,\theta_1+\theta_2,3\theta_1\)
are independent. At a local maximum with nonzero \(F\), these terms
have the same phase. Indeed, after rotating the sum to the positive
axis, the first derivative in each summand's phase forces it to be
real. A negative summand would make the absolute value strictly
increase upon a sufficiently small rotation of that term alone.

If \(q_3=0\), a first-order variation in \(q_3\) aligned with \(F\)
increases its modulus while the normalization costs only second order.
If exactly one of \(q_1,q_2\) is zero, a first-order variation of the
vanishing coefficient has the same effect. If both vanish, scalar
and variable rotations make \(q_3>0\), and the two-coordinate
optimization gives \(q_0^2=a/s\) and \(q_3^2=1/s\). Perturb
\(q_1=q_2=\varepsilon>0\) and rescale \(q_0,q_3\) by
\(\sqrt{1-2\varepsilon^2}\). The coefficient of \(\varepsilon^2\) in the resulting gain is
\[
 s\big(aq_0^{s-2}-F\big)>0,
 \qquad
 \frac{a q_0^{s-2}}{F}=\frac{a}{s q_0q_3}=\sqrt a>1.
\]
This is impossible at a maximizer. Thus all three coefficients are
nonzero; variable rotation makes \(q_1>0\), and phase alignment then
makes \(q_2,q_3>0\).

Equation~\eqref{eq:sphere-lagrange} gives \(\partial F/\partial q_i=sM q_i\).
Put \(x=q_1/q_0\), \(y=q_3/q_0\). The equations for \(q_3,q_2\)
give \(M=q_0^{s-2}/y\) and \(q_2=a x y q_0\), respectively.
The remaining equations are
\[
 a^2y^2+\frac{a(a-1)}2 x y=1,
 \qquad a y^2+a^2(a-1)x^2y^2
       +\frac{a(a-1)(a-2)}6x^3y=1.
\]
Eliminate \(y^2\) using \(u=x/y\). This yields
\eqref{eq:third-cubic} and the stated expressions for \(x,y\).

To identify the root, note \(P_a(0)=-a^2<0\) and
\(P_a(2a)=a^3(a+1)/3>0\). If \(a\geq2\),
Descartes' rule gives exactly one positive root.
If \(1<a<2\), the rule gives at most two positive roots, and
the signs at zero, \(2a\) and infinity give exactly two.
Write
\[
 U=\frac{2(a+1)}{2-a},\qquad
 P_a(U)=\frac{(2a-1)(4a^2-a+4)}{3(a-2)^2}>0 .
\]
Since \(U>2a\) and \(P_a(U)>0\), the larger positive root is greater than \(U\).

The reciprocal of the normalized polynomial is
\(T(z)=z^3+xz^2+a x y z+y\). Its roots are in the open disk,
so \(0<y<1\). With \(T^*(z)=z^3\overline{T(1/\bar z)}\),
Rouch\'{e}'s theorem on \(\T\) shows that \(T-yT^*\) has all
three zeros in the disk, since \(|yT^*|=y|T|<|T|\).
Dividing by \(z\) gives
\[
 (1-y^2)z^2+x(1-a y^2)z+(a-1)xy .
\]
Its root product has modulus less than one, so
\((a-1)xy<1-y^2\). When \(1<a<2\), substituting the equation
\(y^{-2}=a^2+a(a-1)u/2\) shows this is exactly \(u<U\).
The larger root is therefore impossible. Existence of an
extremizer and the uniquely determined positive coefficients now
give the asserted value, its attainment, and all equality cases.
\end{proof}

At \(p=2/3\), one has \(a=2\) and \(u=1+\sqrt{33}/3\).
Direct substitution recovers the Brevig--Saksman value
\[
 C(3,2/3)=\sqrt{\frac{2(1103+33\sqrt{33})}{1153}} .
\]
\section{The sharp fourth coefficient}\label{sec:quartic}

For \(a>1\) and \(0<v<1\), write
\[
 \delta=a-1,\quad A_j=\binom aj,\quad
 k=\frac{\delta v}{2(1-v)},\qquad
 u=\frac{6(1-v)^2(1+v)}
          {v\delta\bigl((a-2)+(2a-1)v\bigr)}.
\]
Whenever the denominator of \(u\) is nonzero, put
\begin{align}
 R_a(v)&=\frac va\bigl[v+2A_2vu+
                   (A_2k^2+3A_3k+A_4)u^2\bigr]-1, \label{eq:quartic-residual}\\
 N_a(v)&=2av\delta\bigl((a-2)+(2a-1)v\bigr)^2 R_a(v).
 \label{eq:quartic-poly}
\end{align}
The expression \(N_a(v)\) extends to a polynomial of degree at most six.

\begin{theorem}\label{thm:quartic}
Let \(0<p<1\) and \(a=2/p-1\). For \(1<a<2\), let \(v\) be the
unique root of \(N_a\) in
\((\sqrt{(2-a)/(2a-1)},1)\). For \(2\le a<3\), let \(v\) be the
unique root in \((1/10,1)\). For \(a\ge3\), let \(v\) be the
unique root in \((0,1)\). At this \(v\), define
\[
 S=1+u+(ku)^2+v^2u+(v/a)^2,\qquad
 Q(z)=S^{-1/2}(1+\sqrt u\,z+ku\,z^2+v\sqrt u\,z^3+(v/a)z^4).
\]
Then the sharp constant for the fourth Taylor coefficient is
\[
 C(4,p)=\frac{a}{v}\,S^{-(a-1)/2}.
\]
The norm-one extremizers are exactly \(\eta Q(\xi z)^{2/p}\), where
\(|\eta|=|\xi|=1\), with the power chosen positive at zero.
\end{theorem}

\begin{proof}
Take a norm-one extremizer and write it as \(Q^{a+1}\), where \(Q\) is
a degree-four polynomial of norm one in \(H^2\), positive at zero,
and zero-free on the closed disk, by Proposition~\ref{thm:boundary}.
Write
\[
 P=Q/Q(0)=1+xz+bz^2+cz^3+dz^4.
\]
If \(x=0\), the first stationary equation gives \(c=0\).
In this case \(Q^{a+1}\) is an analytic function of \(z^2\),
with the same \(H^p\) norm as the corresponding function of \(z\).
The modulus of its fourth coefficient is at most \(C(2,p)\). The strict inequality
\(C(2,p)<C(4,p)\) follows from the outer extremal theorem:
multiplying a second-coefficient extremizer by \(z^2\) preserves
its norm and gives a function that cannot be extremal for the
fourth coefficient. Thus we may rotate the argument to make \(x>0\).

Let \(d=yw\), where \(y>0\) and \(|w|=1\), and write \(v=ay\),
\(u=x^2\). The reflected polynomial identity
\(P^a\equiv P^*/\bar d\pmod{z^5}\) gives
\begin{align}
 c&=axd,\label{eq:quartic-1}\\
 a b+A_2u&=\bar b/\bar d,\label{eq:quartic-2}\\
 ac+2A_2xb+A_3x^3&=x/\bar d,\label{eq:quartic-3}\\
 ad+A_2(2xc+b^2)+3A_3ub+A_4u^2&=1/\bar d.\label{eq:quartic-4}
\end{align}
Also the phase of the fourth coefficient of \(Q^{a+1}\) is \(w\).
From \eqref{eq:quartic-3}, set
\[
 b=Bw-Du,\qquad B=\frac{1/y-a^2y}{a(a-1)},\qquad
 D=\frac{a-2}{6}.
\]
Equation~\eqref{eq:quartic-2} becomes
\[
 (aB+Du/y)w=B/y-\frac{a(2a-1)}6u .
\]
If \(w\) is nonreal, both real quantities on the two sides must
vanish. Since \(u>0\), this gives \(1<a<2\),
\[
 v^2=\frac{2-a}{2a-1},\qquad
 u=\frac{18}{(2a-1)(2-a)}\ge16.
\]
The inequality uses \((2a-1)(2-a)\le9/8\).
Every reciprocal root of \(P\) has modulus strictly less than one,
so \(|x|<4\). This excludes the nonreal case.

Now \(w\in\{-1,1\}\), and \(b,c,d\) are real. The equality \(v=1\)
is impossible for \(w=1\) by \eqref{eq:quartic-2}.
For \(v=1,w=-1\), equations~\eqref{eq:quartic-2} and
\eqref{eq:quartic-3} give both \(b=-A_2u/(2a)\) and
\(b=-A_3u/(2A_2)\). This would require
\(A_3=A_2^2/a\), or \(2(a-2)=3(a-1)\), which is impossible.

We next use local phase variation. Keeping all coefficient moduli
fixed, change the phase \(\theta\) of \(b\) and jointly vary the phases
of \(c,d\). These variations are admissible in a neighborhood,
because the polynomial is zero-free on the closed disk.
If \(b\ne0\), put \(K=|b|/u\). Apart from a common positive factor, the
part of the coefficient of \(Q^{a+1}\) that depends on \(\theta\)
is
\[
 Y(\theta)=L e^{2i\theta}+M e^{i\theta}+H,\qquad
 L=K^2,\quad M=(a-1)K,\quad H=\frac{(a-1)(a-2)}{12}.
\]
The other part, coming from \(c,d\), has phase \(w\) by
\eqref{eq:quartic-1}. In the present real cases \(b\ne0\):
equation~\eqref{eq:quartic-2} with \(b=0\) would give \(A_2u=0\).
At the original phase \(\theta_0\in\{0,\pi\}\), varying the joint
phase of \(c,d\) forces \(wY(\theta_0)\ge0\). If equality holds and
\(Y'(\theta_0)\ne0\), the Hessian of the squared coefficient modulus
has zero diagonal entry in the joint-phase direction and nonzero
mixed entry, contradicting a local maximum. The cases below have
\(Y'(\theta_0)\ne0\). Consequently \(wY(\theta_0)>0\).
One can then vary the joint phase continuously so that the other
part remains aligned with \(Y(\theta)\). Thus \(\theta_0\) is a
local maximum of \(|Y(\theta)|\).

As a function of \(X=\cos\theta\), its squared modulus is
\[
 L^2+M^2+H^2+2LH(2X^2-1)+2M(L+H)X .
\]
The derivatives at \(X=-1\) and \(X=1\), respectively, are
\[
 2K[(a-1)K^2+H(a-1-4K)],\qquad
 2K[(a-1)K^2+H(a-1+4K)].
 \tag{\(\dagger\)}\label{eq:quartic-derivs}
\]
A local phase maximum at \(\theta=\pi\) requires the first derivative
to be nonpositive. A maximum at \(\theta=0\) requires the second
derivative to be nonnegative.

For \(w=-1\), equations~\eqref{eq:quartic-2} and
\eqref{eq:quartic-3} give
\[
 b=-Ku,\qquad K=\frac{(a-1)v}{2(1+v)},\qquad
 u=\frac{6(1-v)(1+v)^2}
         {v(a-1)\bigl((2a-1)v-(a-2)\bigr)} .
\]
The denominator cannot vanish, since the linear equation for \(u\)
has nonzero right side when \(v\ne1\). If \(v>1\), the expression has
\(u<0\), so \(0<v<1\).
We have \(Y'(\pi)=i(2L-M)\ne0\) since \(M/L=2(1+v)/v>4\).
If \(a\ge2\), the first derivative in \eqref{eq:quartic-derivs} is
strictly positive. If \(1<a<2\), its required nonpositivity implies
\(v^2\le(2-a)/(2a-1)\). In this range of \(a\), the displayed
expression for \(u\) strictly decreases on \(0<v<1\), because the
logarithmic derivative is
\[
 -\frac1{1-v}+\frac2{1+v}-\frac1v-
       \frac{2a-1}{(2a-1)v+2-a}<0 .
\]
At \(v=\sqrt{(2-a)/(2a-1)}\), it equals
\(18/((2a-1)(2-a))\ge16\). This again contradicts \(x<4\).

Suppose \(w=1\) and \(v>1\). Then \(b=-Ku\) with
\(K=(a-1)v/(2(v-1))\). In this case
\(Y'(\pi)\ne0\) since \(M/L=2(v-1)/v<2\).
The first derivative in \eqref{eq:quartic-derivs} is positive if
\(a\le2\), by its signs. For \(a>2\), nonpositivity would imply
\(3(a-1)v^2\le(a-2)(v^2-1)\), which is impossible.
It follows that \(w=1\) and \(0<v<1\).
Equations~\eqref{eq:quartic-1}--\eqref{eq:quartic-3} now give the
stated \(k,u\) and \(b=ku>0\). Again the equation for \(u\) has
nonzero right side, so its denominator is nonzero.
At \(\theta_0=0\), the derivative
\(Y'(0)=i(2L+M)\ne0\).
For \(1<a<2\), the second derivative test in
\eqref{eq:quartic-derivs} requires \(v^2\ge(2-a)/(2a-1)\).
At equality \(u=18/((2a-1)(2-a))\ge16\), so the inequality is strict.
The last jet equation is \(R_a(v)=0\).

It remains to count roots in the allowed intervals. Under
\(v=t/(1+t)\), the coefficient list of
\((1+t)^6N_a(t/(1+t))\), in descending order, is
\begin{equation}\label{eq:quartic-table}
 \begin{array}{ll}
 t^6&-18(a-1)^4,\\
 t^5&-6(a-1)^3(11a-7),\\
 t^4&-2(a-1)^2(46a^2-127a+16),\\
 t^3&-2(a-1)(30a^3-169a^2+193a-4),\\
 t^2&-3a(6a^3-61a^2+146a-99),\\
 t^1&-2a(a-2)(a^2-18a+29),\\
 t^0&3a(a-3)(a-2).
 \end{array}
\end{equation}
These identities follow by substitution in
\eqref{eq:quartic-residual}--\eqref{eq:quartic-poly}.
For \(a\ge3\), the \(t^4\) coefficient is negative. As \(a\)
increases from 3, the polynomials in parentheses in the
\(t^3,t^2,t^1\) rows each change sign from negative to positive
at exactly one value. Their sign changes are in \((4,5)\),
\((7,8)\), and \((16,17)\), respectively. To verify uniqueness,
substitute \(a=3+r\) and apply Descartes' rule of signs.
Their descending coefficient lists are
\((30,101,-11,-136)\), \((6,-7,-58,-48)\), and \((1,-12,-16)\).
Thus \eqref{eq:quartic-table} has exactly one sign change
after omitting zeros. At \(a=2\) it also has one sign change,
and \(N_2(1/10)>0\) by the evaluation below.

For \(2<a<3\), the \(t^3,t^2,t^1\) rows are positive.
Indeed, upon substituting \(a=(2+3r)/(1+r)\) and clearing the
denominator, the two cubics in these rows have lists
\((-136,-397,-285,-54)\) and \((-48,-86,-35,-3)\). The quadratic in
the \(t^1\) row is decreasing and negative on \([2,3]\).
Hence there are at most two positive roots in \(t\), counted with
multiplicity. We have \(N_a(0)<0\), \(N_a(1)<0\), and \(N_a(1/10)>0\).
For the last sign, the polynomial
\(10^6(1+r)^4N_{(2+3r)/(1+r)}(1/10)\) is
\[
 8856360r^4+14383161r^3+4718790r^2-315261r+204750 .
\]
Its quadratic tail is positive by its negative discriminant,
so this sign holds also at \(a=2\). It follows that, for \(2<a<3\),
one root lies in \((0,1/10)\) and one in \((1/10,1)\).
The expression \(u(v)\) strictly decreases with \(v\) on \((0,1)\),
as is directly seen by logarithmic differentiation.
Also \(u(1/10)\ge u(1/10)|_{a=3}=891/50>16\) because its
denominator increases with \(a\). The smaller root is inadmissible.

For \(1<a<2\), set \(r=\sqrt{(2-a)/(2a-1)}\), so
\(a=(2+r^2)/(1+2r^2)\). Substitute \(v=(r+t)/(1+t)\).
Multiplication of \(N_a(v)\) by
\((1+t)^6(1+2r^2)^4/[9(1-r)^4]\) yields the following
coefficient list:
\[
\begin{array}{ll}
t^6&-2(1+r)^4,\\
t^5&-2(1+r)^3(5r^2+6r+5),\\
t^4&2(1+r)^2(2r^2-9r+2)(3r^2+r+3),\\
t^3&2(1+r)(6r^6+24r^5-31r^4+52r^3-31r^2+24r+6),\\
t^2&2r^8+36r^7+69r^6+14r^5+174r^4+14r^3+69r^2+36r+2,\\
t^1&2r(1+r)(r^2+8r+1)(2r^4-2r^3+5r^2-2r+2),\\
t^0&2r^2(1+r)^2(2r^4-2r^3+5r^2-2r+2).
\end{array}
\]
The \(t^3,t^2,t^1,t^0\) coefficients are positive. For the two
mixed-sign factors, substitute \(r=h/(1+h)\) and clear the
denominators. The resulting polynomials are
\(50h^6+150h^5+269h^4+288h^3+179h^2+60h+6\) and
\(5h^4+10h^3+11h^2+6h+2\).
The signs therefore have exactly one change, regardless of
the sign of the \(t^4\) coefficient.

An extremizer exists, so in each case the root in the stated
interval is realized by an extremal factor. The equations show
that this root determines all four normalized coefficients.
Finally, \(Q^{a+1}\equiv C(4,p)QQ^*\pmod{z^5}\) at the positive
phase normalization. Comparing constant terms gives
\(C(4,p)=Q(0)^{a-1}/(v/a)\). This proves the formula and the
classification.
\end{proof}
\subsection{A sharp constant not expressible by radicals}

For \(p=1/2\), the sharp constant \(C(4,p)\) is algebraic of degree
five with nonsolvable Galois group. In fact, \(a=3\) gives
\(N_3(v)=8vJ(v)\), where
\[
 J(v)=45v^5-10v^4-49v^3-28v^2-6v+12.
\]
The monic reduction of \(J\) modulo 7 is \(X^5-X^4-2X-3\).
Its gcd with \(X^{49}-X\) is one, as certified by the successive
monic remainders
\[
 X^4-3X^3+3X^2-3X+3,\quad
 X^3-X^2-2X-3,\quad X^2+X-1,\quad X+2,\quad 1 .
\]
A reducible quintic has a factor of degree at most two.
Every monic irreducible over \(\mathbb F_7\) of degree one
or two divides \(X^{49}-X\), so \(J\) is irreducible.
The monic reduction modulo 47 factors as
\[
 (X^2-3X-1)(X^3+8X^2-21X+6).
\]
The quadratic is irreducible since its discriminant is 13 and
\(13^{23}\equiv-1\pmod {47}\).
The cubic is irreducible since the monic Euclidean remainders
starting with \(X^{47}-X\) and that cubic are
\(X^2+8X-14,\ X-21,\ 1\).
The reductions are squarefree and retain their degree. The
factorization-to-cycle theorem gives a 5-cycle and an element
of cycle type \((2)(3)\) in the Galois group of \(J\).
The cube of the latter is a transposition.
Its conjugates under the 5-cycle are the edge transpositions
of a connected graph on five vertices, and hence generate
\(\mathfrak S_5\).

At the specified root of \(J\), the formula in
Theorem~\ref{thm:quartic} gives
\[
 C(4,1/2)=
 \frac{27(5v+1)^2}{135v^6-2v^5-17v^4+64v^3-45v^2+198v+27}.
\]
The denominator reduced modulo \(J\) has coefficient \(1226/9\)
at \(v^4\), whereas the numerator has degree two.
Thus the ratio cannot equal a nonzero rational number, since
such an equality would give a nonzero polynomial of degree
at most four vanishing at \(v\). As \([\mathbb Q(v):\mathbb Q]=5\),
it follows that \(\mathbb Q(C(4,1/2))=\mathbb Q(v)\).
Its normal closure has group \(\mathfrak S_5\), proving the claim.
\section{Exact computation and arithmetic of all coefficient norms}\label{sec:arithmetic}

Fix \(n\geq1\), let \(s=2/p\) and \(a=s-1\), and write
\[
 P(z)=1+r_1z+\cdots+r_nz^n,qquad
 d=P^*(0)=\overline{r_n},\qquad U=\|P\|_2^2 .
\]
Here \(P^*(z)=z^n\overline{P(1/\bar z)}\). Let \(\mathcal S_{n,p}\)
be the set of these coefficient vectors for which \(P\) is zero-free
on \(\overline{\D}\), \(d\ne0\), and
\begin{equation}\label{eq:arithmetic-jet}
 P^a\equiv P^*/d\pmod{z^{n+1}} .
\end{equation}
The coefficients of the left side through degree \(n\) are obtained
by the finite generalized binomial expansion. Consequently
\eqref{eq:arithmetic-jet}, after multiplying by \(d\), is a polynomial
system over the real field generated by \(p\), in the real and
imaginary parts of \(r_1,\ldots,r_n\). The closed-disk root exclusion
is also semialgebraic, by eliminating the two real variables in a
potential zero \(x+iy\) with \(x^2+y^2\leq1\).

\begin{theorem}\label{thm:hardy-arithmetic}
Choose one point \(P_j\) from each cell of any finite semialgebraic
cell decomposition of \(\mathcal S_{n,p}\), and write
\(d_j=P_j^*(0)\), \(U_j=\|P_j\|_2^2\). Then
\[
 C(n,p)=\max_j\frac{1}{|d_j|U_j^{(s-2)/2}} .
\]
Every scalar multiple of modulus one of \((P/\sqrt U)^s\) from a
maximizing cell attains equality, and these exhaust the norm-one
extremizers. If \(p\) is rational, the norm is algebraic. If \(p\)
is algebraic irrational, the norm is transcendental. For every
algebraic \(p\in(0,1)\), the maximizing cells and a finite exact
expression for the norm can be found by an algorithm.
\end{theorem}

\begin{proof}
Theorem~\ref{thm:boundary} and the coefficient stationarity
give \(\mathcal S_{n,p}\ne\varnothing\) and place every
normalized extremizer in it. We verify the required converse at
the stationary level. Let \(P\in\mathcal S_{n,p}\), put
\(Q=P/\sqrt U\), \(f=Q^s\), and write \(L(f)=[z^n]f\).
Multiplying \eqref{eq:arithmetic-jet} by \(P\) gives
\[
 f\equiv \frac{U^{1-s/2}}{d}\,QQ^*\pmod{z^{n+1}} .
\]
Since the coefficient at degree \(n\) in \(QQ^*\) is \(1\),
we have \(L(f)=U^{1-s/2}/d\ne0\).
On \(\T\) we have \(z^{-n}QQ^*=|Q|^2\). It follows that for every
function \(g\) holomorphic on a neighborhood of the closed disk,
\[
 L(fg)=L(f)\int_{\T}|Q|^2g\,dm .
\]
Along a continuously differentiable path in \(\mathcal S_{n,p}\),
use \(g=s\dot Q/Q\). The derivative of the logarithm of the modulus
of the coefficient is
\[
 \frac{d}{dt}\log|L(Q^s)|
 = s\Re\int_{\T}\overline Q\dot Q\,dm
 =\tfrac{s}{2}\frac{d}{dt}\|Q\|_2^2=0 .
\]
Semialgebraic cells admit piecewise continuously differentiable
paths between their points \cite{BCR}. Thus the displayed value in the theorem
is constant on each cell. Each sample is admissible in the original
Hardy problem, and some cell contains a global extremizer. This
proves the exact max formula and the equality statement.

If \(p\) is algebraic, the defining variables of a cylindrical
algebraic decomposition and its sample points may all be taken
algebraic, with real and imaginary parts represented by real-root
isolations. Since \(d_j\ne0\), every \(U_j>1\).
If \(p\) is rational then \((s-2)/2\) is rational and all candidate
values are algebraic. If \(p\) is algebraic irrational then
\((s-2)/2\) is algebraic irrational, and Gelfond--Schneider shows
that each \(U_j^{(s-2)/2}\) is transcendental. Multiplication and
division by nonzero algebraic numbers preserve transcendence, so
each candidate value is transcendental.

For algebraic irrational \(p\), comparing two positive candidate
values reduces to determining whether \(A B^{(s-2)/2}\) is larger
than, equal to, or smaller than \(1\), with \(A,B>0\) algebraic.
If \(B=1\), algebraic comparison decides. If \(B\ne1\), the
Gelfond--Schneider theorem excludes equality, and certified
approximations to the logarithms and the exponent decide the strict
comparison after finitely many refinements. Together with exact
algebraic comparison for rational \(p\), this proves the algorithmic
claim.
\end{proof}

For an arbitrary real \(p\in(0,1)\) the formula is a finite exact
representation. The termination claim about comparisons is made
for algebraic \(p\). This resolves uniform exact computation in that
arithmetic range and gives individual finite representations in
every degree. The closed formulas in low degrees give explicit
information beyond the general elimination procedure.

\subsection{Finite linear combinations of coefficients}

\begin{corollary}\label{cor:hardy-arithmetic-general}
Let \(L_\gamma(f)=\sum_{k=0}^n\gamma_k[z^k]f\) with algebraic
complex \(\gamma_k\), \(n\geq1\) and \(\gamma_n\ne0\).
For \(0<p<1\), the norm of this functional on \(H^p\) is
algebraic for rational \(p\) and transcendental for algebraic
irrational \(p\). In both ranges it admits a finite exact
computation of the type above.
\end{corollary}

\begin{proof}
The inner--outer concavity argument and the polynomial perturbation
argument work for \(L_\gamma\). They give an outer maximizing
function \(f\) of norm one for which \(|f|^p\) is a nonnegative
trigonometric polynomial of degree at most \(n\). Indeed, its
moments satisfy \(L_\gamma(z^j f)=L_\gamma(f)\int|f|^p z^j\,dm\)
for all \(j\geq0\), with the left side zero for \(j>n\).
At \(j=n\), it is \(\gamma_n f(0)\ne0\). Thus the trigonometric
degree is \(n\), and spectral factorization and outerness give
\(f=\eta Q^s\) with \(Q=P/\sqrt U\) as above. For this functional
we allow roots of \(P\) on \(\T\).

Put \(r_0=1\),
\[
 L_j=L_\gamma(z^jP^s),\qquad
 c_j=\sum_{k=0}^{n-j}r_k\overline{r_{k+j}},\qquad 1\leq j\leq n .
\]
The polynomial stationarity system is \(UL_j=L_0c_j\) for
\(1\leq j\leq n\), together with \(L_0\ne0\) and root exclusion
in the open disk. Its solutions include all extremizers and satisfy
\(r_n\ne0\) because \(U\gamma_n=L_0\overline{r_n}\).
For any solution, the same weighted integral identity as in the
coefficient proof follows for \(g\) holomorphic on a neighborhood
of \(\overline{\D}\), with \(L=L_\gamma\) and \(f=Q^s\).

For a continuously differentiable path of solutions, replace
\(s\dot Q/Q\) temporarily by
\(g_\rho(z)=s\dot Q(\rho z)/Q(\rho z)\), \(0<\rho<1\).
These functions are holomorphic past the boundary.
For each factor \(1-w\) of the boundary value of \(P\) we have
\(|w|\leq1\) and
\[
 |1-\rho w|^2-\rho|1-w|^2
   =(1-\rho)(1-\rho|w|^2)\geq0 .
\]
Hence \(|Q(\zeta)|/|Q(\rho\zeta)|\leq\rho^{-n/2}\) on \(\T\).
This bounds \(|Q|^2|g_\rho|\) by
\(s\rho^{-n/2}|Q|\max_{|z|\leq1}|\dot Q(z)|\).
Dominated convergence on \(\T\) and convergence of Taylor jets
inside the disk give
\[
 \frac{d}{dt}\log|L_\gamma(Q^s)|
   =s\Re\int_{\T}\overline Q\dot Q\,dm=0 .
\]
Cell decomposition therefore yields finitely many exact candidate
values \(|L_0|U^{-s/2}\). Their numerators are nonzero algebraic
at algebraic samples, and \(U>1\) by \(r_n\ne0\).
The same Gelfond--Schneider argument and logarithm comparison
complete the proof.
\end{proof}

\begin{corollary}\label{cor:hardy-phase-finite}
For fixed \(n\), the function \(p\mapsto C(n,p)\) is real analytic
on \((0,1)\) outside a finite set.
\end{corollary}

\begin{proof}
Apply cylindrical algebraic decomposition to the jet system with
the parameter \(s\) ordered first. On finitely many open parameter
intervals, choose one semialgebraic section of each nonempty cell
fiber. After subdividing each interval finitely, these sections,
their positive \(U\) values and their positive \(|d|\) values
are real analytic and algebraic over \(\mathbb R(s)\).
By the theorem the coefficient norm is the maximum of the associated
positive candidate functions. The logarithm of the ratio of two
of them has the form
\(h(s)=\log A(s)+(s-2)\log B(s)/2\)
for positive real-analytic algebraic functions \(A,B\).
Its second derivative is semialgebraic: differentiating twice
removes the logarithms. If \(h''\) is not identically zero, it
has finitely many zeros on the interval by semialgebraicity and
analyticity. Rolle's theorem implies that \(h\) has finitely many
zeros there. If \(h''\equiv0\), then \(h\) is affine, so the two
branches tie identically or cross at most once. Subdivide at all
remaining crossings to obtain the claimed analyticity.
\end{proof}
\section{Limits at the endpoints of the exponent range}

\begin{proposition}\label{prop:smallp}
For each fixed integer \(n\ge1\),
\[
 C(n,p)=\frac{1}{n!}\left(\frac{2n}{ep}\right)^{n/2}
            \bigl(1+O_n(\sqrt p)\bigr)
 \qquad(p\downarrow0).
\]
\end{proposition}

\begin{proof}
Write \(Q= (1+\sum_{j=1}^n r_jz^j)/\sqrt{1+R^2}\), with
\(R^2=\sum_{j=1}^n |r_j|^2\), for the factor of an extremizer.
In the binomial expansion for \(\coef n Q^s\), the term involving
\(j\) factors from \(\sum_{k=1}^n r_kz^k\) is bounded in absolute
value by \(O_n(s^j R^j (1+R^2)^{-s/2})\).
For \(s>j\), the maximum of \(R^j(1+R^2)^{-s/2}\) occurs at
\(R^2=j/(s-j)\), so each term with \(j<n\) contributes
\(O_n(s^{(n-1)/2})\) in total. The term with \(n\) factors is
\(\binom sn r_1^n(1+R^2)^{-s/2}\), and is bounded by
\[
 \binom sn\left(\frac{n}{s-n}\right)^{n/2}
     \left(1+\frac{n}{s-n}\right)^{-s/2}.
\]
The degree-one polynomial with \(r_1=\sqrt{n/(s-n)}\) and
\(r_2=\cdots=r_n=0\) gives this exact value for its \(n\)-th
coefficient, and its power is an admissible comparison function
for \(s>2n\). Expanding the displayed expression at \(s=\infty\)
gives the result.
\end{proof}

For a density \(W\) on \(\T\), use the convention
\(\Ent(W)=\int_{\T}W\log W\,dm\), with \(0\log0=0\).
Let \(\mathcal U_n\) be the degree-\(n\), \(H^2\)-normalized
polynomials whose roots are all on \(\T\).

\begin{theorem}\label{thm:entropy}
For each \(n\ge1\),
\[
 \lim_{p\uparrow1}\frac{\log C(n,p)}{1/p-1}
      =\max_{Q\in\mathcal U_n}\Ent(|Q|^2).
\]
Every coefficientwise limit of the normalized extremal polynomial
factors as \(p\uparrow1\) attains the maximum on the right.
\end{theorem}

\begin{proof}
Every \(Q\in\mathcal U_n\) is self-reciprocal up to phase.
Consequently \(|\coef n Q^2|=1\). Use \(Q^2\), divided by
its \(H^p\) quasi-norm, as a comparison function in \eqref{eq:def}.
For \(W=|Q|^2\), the result is
\[
 \log C(n,p)\ge -\frac1p\log\int_{\T} W^p\,dm.
\]
A Taylor expansion in \(p\) at 1 gives
\((1/p-1)\Ent(W)+O_n((1-p)^2)\) on the right. The error
is uniform for \(0\le W\le n+1\), \(\int W\,dm=1\), since
\(x^t(\log x)^2\) is bounded on \([0,n+1]\) when \(t\) is
in a fixed neighborhood of 1. The set of unit-root polynomials
with norm one is compact, and entropy is continuous on it.

For the upper bound, take an extremal factor \(Q_p\) and put
\(W_p=|Q_p|^2\). Then
\[
 C(n,p)\le \|Q_p^{2/p}\|_1
       =\int_{\T} W_p^{1/p}\,dm
       =1+(1/p-1)\Ent(W_p)+O_n((1-p)^2).
\]
The bound \(W_p\le n+1\) guarantees the same uniform expansion.
The normalized polynomials have convergent subsequences in
coefficient space. Their constant terms are at least \(2^{-n}\):
the constant-term-one polynomial is a product of \(n\) factors
\(1-\alpha z\), \(|\alpha|\le1\), and has sup norm at most \(2^n\).
Every limiting \(Q_0\) is zero-free in \(\D\) by Hurwitz.
The two inequalities show \(C(n,p)\to1\).
Local convergence of analytic powers on smaller disks implies
\(|\coef n Q_0^2|=1\). Since \(\|Q_0\|_2=1\), equality in
Cauchy--Schwarz for the coefficient \([z^n]Q_0^2\) gives
\(Q_0^{*n}=\xi Q_0\) for some \(|\xi|=1\).
Its constant term is nonzero, so \(Q_0\) has degree \(n\).
Reciprocal pairing and disk zero-freeness put all zeros on \(\T\).
Continuity of entropy then gives the upper bound along any
subsequence and proves the claim about limits.
\end{proof}

This maximum differs from the unit-root entropy minimum in the
Agler--McCarthy program for bounded zero-free coefficients,
proved by Lei and Zhang \cite{LeiZhang}.

\paragraph{Evaluation in degree three.}
As \(a\downarrow1\), the selected root \(u\) in Theorem~\ref{thm:third}
converges to the unique root \(u_*\in(1,2)\) of
\(u^3-6u^2+3u+6=0\) on \([0,2]\).
Indeed \(0<u<2a\), and the limiting cubic has one negative root,
one root in \((1,2)\), and one root in \((5,6)\).
The limit polynomial is
\((1+u_*z+u_*z^2+z^3)/\sqrt{2+2u_*^2}\).
With \(\delta=a-1\), the formula for \(y\) gives
\(-\log y=\delta(1+u_*/4)+o(\delta)\).
Since \(1/p-1=\delta/2\), the entropy maximum from
Theorem~\ref{thm:entropy} is
\[
 2+u_*/2-\log(2+2u_*^2).
\]

\paragraph{Evaluation in degree four.}
For the fourth coefficient, the limit in
\(\log C(4,p)/(1/p-1)\) can be evaluated in terms of the unique
root \(h_0\in(0,1)\) of
\[
 12h_0^4-50h_0^3+65h_0^2-12h_0-9=0.
\]
Set \(u_0=12h_0^2/(3-h_0)\) and \(S_0=2+2u_0+(u_0/(2h_0))^2\).
Then the limit is \(2+2h_0-\log S_0\), and the normalized limiting
fourth-coefficient factor is
\[
 S_0^{-1/2}
 \left(1+\sqrt{u_0}\,z+\frac{u_0}{2h_0}z^2+\sqrt{u_0}\,z^3+z^4\right).
\]
Indeed, writing \(e=a-1\) and \(v=1-eh\) in Theorem~\ref{thm:quartic},
the lower bound on \(v\) gives \(0<h<3/2+o(1)\).
Direct substitution in \eqref{eq:quartic-poly} gives
\[
 e^{-4}N_{1+e}(1-eh)=
 2(12h^4-50h^3+65h^2-12h-9)+O(e),
\]
uniformly for bounded \(h\). The polynomial in parentheses has the
stated single root on \([0,3/2]\). For \((0,1)\), the substitution
\(h=t/(1+t)\) gives descending coefficients
\((6,8,-25,-48,-9)\) after clearing denominators. For \([1,2]\)
the substitution \(h=(1+2t)/(1+t)\) gives
\((19,44,71,40,6)\), all positive.
Thus \(h\to h_0\). The formulas in Theorem~\ref{thm:quartic} give
\(u\to u_0\), \(k\to1/(2h_0)\), \(S\to S_0\), and
\[
 \log C(4,p)=-\tfrac e2\log S-\log(1-eh)+\log(1+e).
\]
Since \(1/p-1=e/2\), this yields the claimed value.

\section{Scope and relation to prior work}

The third- and fourth-coefficient results use the original supremum
\eqref{eq:def}, with no restriction on the functions in its unit
ball. The proof uses the outer reduction of Brevig and S{\o}vik
\cite{BrevigSovik} and adds a boundary estimate and a phase
exclusion. In particular the all-exponent third formula agrees
at \(p=2/3\) with the value of Brevig and Saksman \cite{BrevigSaksman}.
The quartic analysis selects the unique branch responsible for the
maximum at each exponent. The higher-degree result gives a finite
calculation and an arithmetic classification at algebraic
exponents. Closed formulas comparable to the cubic and sextic
remain to be determined in higher degrees.

The literature comparison covers the original coefficient paper,
the subsequent solution of outerness, the 2026 account of the
sharp coefficient problem, and the VMRC snapshot listed below.
No full-exponent formula for \(C(3,p)\) or \(C(4,p)\), or
all-degree arithmetic dichotomy, was identified in those sources
or in targeted searches. This comparison is limited to the sources
reviewed.

\small
\begin{thebibliography}{99}
\setlength{\itemsep}{2pt}

\bibitem{BCR}
J.~Bochnak, M.~Coste and M.-F.~Roy, \emph{Real Algebraic Geometry},
Springer, 1998.
\href{https://doi.org/10.1007/978-3-662-03718-8}{DOI: 10.1007/978-3-662-03718-8}.

\bibitem{BrevigSaksman}
O.~F.~Brevig and E.~Saksman, Coefficient estimates for \(H^p\)
spaces with \(0<p<1\), \emph{Proc. Amer. Math. Soc.} \textbf{148}
(2020), 3911--3924.
\href{https://doi.org/10.1090/proc/14995}{DOI: 10.1090/proc/14995}.

\bibitem{BrevigSovik}
O.~F.~Brevig and \O.~S{\o}vik, Extremals for linear functionals on
\(H^p\) are outer when \(0<p<1\),
\href{https://arxiv.org/abs/2610.21788}{arXiv:2610.21788}, 2026.

\bibitem{ChenEtAl}
S.~Chen, M.~Huang, L.~Jin and Q.~Li, Sharp constant problems for
harmonic mappings, \href{https://arxiv.org/abs/2609.33159}{arXiv:2609.33159}, 2026.

\bibitem{Duren}
P.~L.~Duren, \emph{Theory of \(H^p\) Spaces}, Academic Press, 1970.

\bibitem{Lang}
S.~Lang, \emph{Algebra}, revised third edition, Springer, 2002.
Appendix 1, Gelfond--Schneider theorem.

\bibitem{LeiZhang}
J.~Lei and T.~Zhang, Proof of the Agler--McCarthy entropy conjecture,
\href{https://arxiv.org/abs/2605.03949v2}{arXiv:2605.03949v2}, 2026.

\bibitem{VMRC248}
Verifiable Mathematical Reasoning Corpus, manuscript [248],
\emph{Sharp third coefficients for zero-free Hardy functions},
\href{https://github.com/QEEDDD/Verifiable-Mathematical-Reasoning-Corpus/tree/e73bae9fd0c319129b19aacc5c317f7abf2865ef}{snapshot e73bae9f},
2026.

\bibitem{VMRC249}
Verifiable Mathematical Reasoning Corpus, manuscript [249],
\emph{Collected Krzy\.{z} manuscripts}, snapshot e73bae9f, 2026.

\end{thebibliography}
\end{document}
